% En-tête LaTeX de la version PDF (moteur : LuaLaTeX).

% Polices de secours pour les caractères absents des polices principales :
% le coréen et les emoji de igposts.csv et du texte (polices de macOS).
\directlua{luaotfload.add_fallback("secours", {
  "AppleSDGothicNeo:mode=harf;",
  "Apple Color Emoji:mode=harf;"
})}
\setmainfont{Latin Modern Roman}[RawFeature={fallback=secours}]
\setsansfont{Latin Modern Sans}[RawFeature={fallback=secours}]
\setmonofont{Menlo}[Scale=MatchLowercase, RawFeature={fallback=secours}]

% Code et sorties : un peu plus petits, et les lignes trop longues passent à
% la ligne au lieu de déborder de la page.
\usepackage{fvextra}
\fvset{fontsize=\small, breaklines=true, breakanywhere=true}
\DefineVerbatimEnvironment{Highlighting}{Verbatim}{commandchars=\\\{\},fontsize=\small,breaklines,breakanywhere}
% Sorties de R (environnement verbatim) : redéfini au début du document,
% après les paquets chargés par Quarto.
\AtBeginDocument{%
  \renewenvironment{verbatim}%
    {\VerbatimEnvironment\begin{Verbatim}[fontsize=\small,breaklines,breakanywhere]}%
    {\end{Verbatim}}}

% Les URL longues peuvent être coupées.
\usepackage{xurl}

% Page de titre : affiliations sous le nom de l'auteur (KOMA-Script).
\publishers{\normalsize Université de Fribourg (2023–2025)\\
  Freie Universität Berlin}

% Point médian de l'écriture inclusive (étudiant·es) : le glyphe de Latin
% Modern est entouré de grands blancs ; on les resserre.
\usepackage{newunicodechar}
\newunicodechar{·}{\kern-.12em\symbol{"00B7}\kern-.12em}

% Plus de souplesse pour couper les lignes avant un long code en ligne.
\setlength{\emergencystretch}{3em}

% Micro-typographie : aide à éviter les lignes qui débordent.
\usepackage{microtype}
